\section{Studi Kasus Manipulasi Bit}
\label{sec:manipulasi-bit}

Tiga persoalan berikut memperlihatkan bagaimana operasi logika,
pergeseran, dan rotasi dipakai bersama untuk menyelesaikan masalah
nyata yang tidak dapat dikerjakan oleh instruksi aritmatika biasa.

\subsection{Memeriksa Paritas Sebuah Byte}

\term{Paritas} menyatakan apakah cacah bit 1 sebuah byte genap atau
ganjil. Pemeriksaan paritas dipakai pada transmisi data untuk mendeteksi
satu bit yang berubah karena gangguan; prosesor 8086 menyediakan
\flag{PF} untuk keperluan ini.

Cara pertama: jalankan \lstasm|TEST AL, AL| agar \eng{flag} disetel
menurut isi \reg{AL} tanpa mengubahnya, lalu periksa \flag{PF}.

\begin{lstlisting}[language=MASM, numbers=none]
    MOV AL, nilai       ; byte yang diperiksa
    TEST AL, AL         ; setel PF menurut cacah bit 1 pada AL
    JP  CETAK_GENAP     ; PF = 1 -> cacah bit 1 genap
    ; jatuh ke sini bila cacah bit 1 ganjil
\end{lstlisting}

Cara kedua tidak bergantung pada \flag{PF}: lipat (\eng{fold}) byte
menjadi setengah lebarnya dengan \asm{XOR}, lalu ulangi sampai tersisa
satu bit.

\begin{lstlisting}[language=MASM, numbers=none]
    MOV AL, nilai       ; byte yang diperiksa
    MOV AH, AL
    SHR AH, 4
    XOR AL, AH          ; lipat nibble tinggi ke nibble rendah
    MOV AH, AL
    SHR AH, 2
    XOR AL, AH
    MOV AH, AL
    SHR AH, 1
    XOR AL, AH
    AND AL, 1           ; AL = 1 bila cacah bit 1 ganjil, 0 bila genap
\end{lstlisting}

Cara ketiga dipakai pada program \file{paritas.asm} di Bagian~Contoh
Program: geser byte satu bit demi satu bit dengan \asm{SHR}, lalu
jumlahkan bit yang keluar dari \flag{CF} memakai \asm{ADC}. Cara ini
paling gamblang karena setiap langkahnya tampak langsung, meskipun
biaya waktunya lebih besar. Perbandingan ketiga cara dirangkum pada
Tabel~\ref{tab:paritas}.

\begin{table}[H]
    \centering
    \caption{Tiga cara pemeriksaan paritas sebuah byte.}
    \label{tab:paritas}
    \small
    \setlength{\tabcolsep}{5pt}
    \begin{tabular}{@{}p{2.6cm}p{4.2cm}p{6.2cm}@{}}
        \toprule
        \textbf{Cara} & \textbf{Instruksi kunci} & \textbf{Catatan} \\
        \midrule
        Flag \flag{PF} & \lstasm|TEST AL, AL| lalu \lstasm|JP| &
        Paling singkat; hanya berlaku untuk byte rendah hasil operasi \\
        \addlinespace
        Pelipatan XOR & \lstasm|XOR AL, AH| tiga kali &
        Tidak bergantung pada \flag{PF}; mudah diperluas ke word dan
        dword \\
        \addlinespace
        Pergeseran dan \asm{ADC} & \lstasm|SHR AL, 1| lalu
        \lstasm|ADC BL, 0| & Paling gamblang; jumlah iterasi sama dengan
        cacah bit operand \\
        \bottomrule
    \end{tabular}
\end{table}

\subsection{\eng{Packing} dan \eng{Unpacking} Dua Digit BCD}

\term{BCD} (\eng{Binary Coded Decimal}) menyimpan bilangan desimal
dengan empat bit per digit. Satu byte menampung dua digit: puluhan pada
nibble tinggi dan satuan pada nibble rendah, yang disebut
\eng{packed BCD}.

Menyusun dua digit menjadi satu byte disebut \eng{packing}, sedangkan
memisahkannya kembali disebut \eng{unpacking}. Keduanya hanya
memerlukan pergeseran dan operasi logika, seperti tampak pada
Tabel~\ref{tab:bcd-packing}.

\begin{table}[H]
    \centering
    \caption{\eng{Packing} dan \eng{unpacking} BCD dengan contoh
    puluhan $= 7$ dan satuan $= 5$.}
    \label{tab:bcd-packing}
    \small
    \setlength{\tabcolsep}{5pt}
    \begin{tabular}{@{}p{3.4cm}p{6.0cm}p{4.0cm}@{}}
        \toprule
        \textbf{Operasi} & \textbf{Instruksi} & \textbf{Hasil} \\
        \midrule
        \eng{Packing} & \lstasm|MOV AL, puluhan| lalu
        \lstasm|SHL AL, CL| (CL = 4) lalu \lstasm|OR AL, satuan| &
        \bin{0111 0101} = \hex{75} \\
        \addlinespace
        Ambil satuan & \lstasm|AND AL, 0Fh| &
        \bin{0000 0101} = \hex{05} \\
        \addlinespace
        Ambil puluhan & \lstasm|SHR AL, CL| dengan CL = 4 &
        \bin{0000 0111} = \hex{07} \\
        \bottomrule
    \end{tabular}
\end{table}

Bentuk programnya adalah sebagai berikut.

\begin{lstlisting}[language=MASM, numbers=none]
    ; --- packing: susun 75 menjadi satu byte BCD ---
    MOV AL, 7           ; digit puluhan
    MOV CL, 4
    SHL AL, CL          ; AL = 70h (0111 0000b)
    OR  AL, 5           ; AL = 75h (0111 0101b) -> 75 dalam BCD

    ; --- unpacking: pisahkan kembali ---
    MOV AH, AL
    AND AH, 0Fh         ; AH = 05h (digit satuan)
    SHR AL, CL          ; AL = 07h (digit puluhan)
\end{lstlisting}

Penyesuaian BCD setelah operasi aritmatika dikerjakan oleh \asm{DAA} dan
\asm{AAM} pada Bab~\ref{chap:aritmatika}. Bentuk BCD juga menguntungkan
saat dicetak: byte \hex{75} tampil sebagai ``75''.

\subsection{\eng{XOR Cipher} Sederhana}

Bila sebuah nilai $A$ di-XOR dengan kunci $K$, hasilnya $B = A \oplus K$.
Nilai $A$ dapat dipulihkan dengan meng-XOR-kan $B$ sekali lagi dengan
kunci yang sama:

\begin{equation}
B \oplus K = (A \oplus K) \oplus K = A \oplus (K \oplus K) = A
\oplus 0 = A
\label{eq:xor-balik}
\end{equation}

Sifat pada Persamaan~\eqref{eq:xor-balik} berarti satu operasi yang sama
dapat dipakai untuk enkripsi dan dekripsi. Inilah bentuk penyandian data
paling sederhana, yang disebut \eng{XOR cipher}.

\begin{studikasusbox}
Sebuah perusahaan menyandikan berkas konfigurasi berisi kata sandi
dengan meng-XOR setiap byte terhadap sebuah kunci tetap. Saat berkas
dibaca kembali, setiap byte di-XOR dengan kunci yang sama sehingga isi
aslinya pulih. Cara ini menyembunyikan teks dari pengamatan sekilas,
tetapi \emph{bukan} pengamanan yang sungguh-sungguh: bila sebagian teks
dapat ditebak, kuncinya terhitung dengan satu operasi \asm{XOR}. Untuk
keperluan nyata diperlukan kunci beruntun (\eng{key stream}) sepanjang
pesannya.
\end{studikasusbox}

Program lengkapnya disajikan pada Bagian~Contoh Program sebagai
\file{xor-cipher.asm}. Kerangka intinya adalah perulangan yang
meng-XOR setiap byte sampai ditemukan penanda akhir string.

\begin{lstlisting}[language=MASM, numbers=none]
    LEA SI, pesan       ; SI menunjuk awal string
ULANG:
    MOV AL, [SI]        ; ambil satu byte
    CMP AL, '$'         ; penanda akhir string tercapai?
    JE  SELESAI
    XOR AL, kunci       ; sandikan byte dengan kunci
    MOV [SI], AL        ; tulis kembali di tempat yang sama
    INC SI
    JMP ULANG
SELESAI:
\end{lstlisting}

\begin{infobox}
Operasi \asm{XOR} punya kebiasaan khas pada kode hasil kompilator:
\lstasm|XOR AX, AX| dipakai untuk mengosongkan \reg{AX}. Instruksi ini
hanya dua byte dan selalu menyetel \flag{ZF}, sedangkan
\lstasm|MOV AX, 0| memerlukan tiga byte. Pola ini sering muncul pada
hasil \eng{disassembly} (Bab~\ref{chap:disassembly}) dan kadang
membingungkan pembaca yang mengira ada operasi logika yang sedang
dilakukan.
\end{infobox}
